\documentclass[11pt]{article}
\usepackage{siunitx}
\usepackage{amsmath}
\usepackage{amssymb}
\usepackage{amsfonts}
\usepackage{subfigure}
\usepackage[x11names]{xcolor}
\usepackage[left=3.5cm,right=3.5cm]{geometry}
\usepackage{framed}
\usepackage{graphicx}
\usepackage{hyperref}
\usepackage{amsthm}
\usepackage{times}
\usepackage{multicol}
\usepackage{booktabs}
\usepackage{tabularx}
\usepackage{multirow}
\usepackage[para,online,flushleft]{threeparttable}
\usepackage{float}
\usepackage[font=footnotesize,labelfont=bf]{caption}

\newtheorem{remark}{Remark}
\colorlet{shadecolor}{Thistle3!50}
\newcommand{\change}[1]{{\color{blue}#1}}
\newcommand{\todo}[1]{{\color{red}#1}}
\graphicspath{{figures/}}
\setlength{\parskip}{0.25em}

\begin{document}
\title{Author Response for Submission\\
``PushAround: Collaborative Path Clearing via Physics-Informed Hybrid Search''}
\author{Author Names Omitted for Anonymous Review.}
\date{}
\maketitle

%========================================
\section*{To Associate Editor}
%========================================
\begin{shaded*}
\noindent \textbf{Q:}
{\itshape
The reviewers acknowledge the contribution of the paper, however, they
have some concerns that need to be addressed.
}
\end{shaded*}

\noindent \textbf{A:}
We sincerely thank the Associate Editor for the careful handling of the
submission and for summarizing the reviewers' constructive comments. The revised
draft addresses every point raised by Reviewers~3 and~4 through clarification of
the WCCG connectivity interpretation, the role and guarantees of the hybrid
search objective, the hardware execution interface, state-based online
replanning, and the comparison protocol. It also introduces dedicated validation
for physics-model mismatch, WCCG connectivity, density statistics, and scaling
with the number of robots. Numerical quantities that require new experiments are
marked in red with \verb|\todo{}| in this initial response draft.

The principal revisions are summarized below.
\begin{itemize}
\item \textbf{Physics-model robustness.} A new planner--execution mismatch study
perturbs mass, friction, robot force limits, and contact dynamics while keeping
the planning model nominal. The main manuscript and Supplementary Material now
contain the protocol and result placeholders. \todo{Insert final perturbation
levels, trial count, and numerical conclusions.}

\item \textbf{WCCG connectivity justification.} Sec.~III-A now interprets WCCG
through $W/2$ configuration-space inflation and explicitly discusses non-convex
boundary decomposition. A new connectivity-validation experiment compares WCCG
with an independent rasterized configuration-space flood-fill reference.
\todo{Insert final agreement/false-positive/false-negative statistics.}

\item \textbf{Search objective and guarantees.} Eq.~(5) is reformulated as a
clearance cost used to guide the search. The revised text states explicitly that
PIHS uses sampled actions, bounded lookahead, a heuristic priority, and
first-feasible termination; hence no global optimality or completeness claim is
made for the underlying continuous hybrid problem.

\item \textbf{Hardware force/control interface and coupled contacts.} The
revision explains that optimized contact forces are feasibility variables rather
than direct hardware setpoints. The Mecanum robots execute planar velocity
commands while tracking mode-defined contact poses. Short-horizon simulation
propagates all movable bodies, so contact-induced motion of neighboring objects
is retained in the successor state.

\item \textbf{Execution discrepancy and replanning.} The previous phrase
``deviation from simulation'' is replaced by explicit state-based monitoring of
pre-contact reachability, contact tracking/stall, task predicates, and updated
WCCG/schedule validity. Motion capture provides the required robot/object states
in the hardware experiments.

\item \textbf{Experimental protocol and statistics.} The manuscript now states
shared reachability/collision checks and computation budgets. The density suite
correctly records 16 workers, an 800~s timeout, and 20,000 control steps. The
revised density experiment increases the number of independent seeds and reports
dispersion. \todo{Insert final seed count and regenerated statistics/figure.}

\item \textbf{Robot-team scaling.} The complexity discussion now explains the
$N$-dependence of contact replacement and the $2N$-dimensional force LP, and a
new simulation table varies $N=1,2,3,4$. \todo{Insert final robot-scaling data.}
\end{itemize}

The detailed responses below provide the corresponding manuscript locations and
revised text in blue. Added experimental values and conclusions that remain to be
collected are shown in red.

\newpage
%========================================
\section*{To Reviewer 3}
%========================================
\begin{shaded*}
\noindent \textbf{Q:}
{\itshape
Overall, this is a good contribution to the field and a well-written
paper. The combination of geometric reasoning, recursive pushing, and
physics-informed search is interesting, and the real-world validation
strengthens the contribution.
}
\end{shaded*}

\noindent \textbf{A:}
We sincerely thank the reviewer for the positive assessment of PushAround and,
in particular, for recognizing the end-to-end integration, recursive pushing,
dense-clutter scalability, and hardware validation. The revision focuses on the
reviewer's two stated weaknesses and four follow-up questions without changing
the central planning framework.

%==============================
\begin{shaded*}
\noindent \textbf{Q:}
{\itshape
Physics-model robustness is insufficiently characterized: The method
relies heavily on simulation-based feasibility validation, but
sensitivity to uncertainty in mass, friction, force limits, contact
dynamics etc. is not systematically evaluated.
}
\end{shaded*}

\noindent \textbf{A:}
Thank you for identifying this important limitation. We agree that the original
hardware experiments demonstrate closed-loop recovery from execution error but
do not isolate sensitivity to individual physics-model parameters. The revised
draft therefore adds a dedicated planner--execution mismatch experiment.

\begin{itemize}
\item \textbf{Mismatch rather than retuning.} Planning always uses the nominal
mass, friction, force-limit, and contact model. Only the execution simulator is
perturbed, so the test measures sensitivity to an incorrect planning model.

\item \textbf{Perturbation families.} Obstacle mass, robot--object friction,
available pushing force, and contact dynamics are varied separately using the
same scene seeds. \todo{Insert final perturbation levels and number of trials.}

\item \textbf{Metrics.} The revision reports success rate, execution time,
number of replans, and push count. This distinguishes outright failures from
cases where online replanning absorbs moderate mismatch.

\item \textbf{Claim scope.} The manuscript no longer relies on the hardware runs
alone to imply general model robustness; it explicitly states that simulation
validation is model dependent and quantifies the sensitivity.
\end{itemize}

\noindent\textbf{Changes in the revised manuscript:}
\begin{itemize}
\item The following statement was added in Sec.~I-B to make the model-dependence
explicit:
\change{Because this physical validation is model dependent, the revised
experiments explicitly quantify sensitivity to planner--execution mismatch in
mass, friction, available force, and contact dynamics.}

\item A new subsection, Sec.~IV-A5, was added with the following protocol:
\change{Because PIHS validates candidates with a physics simulator, robustness
is evaluated under deliberate planner--execution model mismatch. The planner
retains the nominal model used to generate the schedule, while the execution
simulator perturbs one parameter family at a time: obstacle mass, robot--object
friction, available robot force, and contact dynamics. This protocol
distinguishes sensitivity to model error from retuning the planner to the
perturbed model.}

\item A result figure and the corresponding discussion have been placed in the
main paper, and the full table/protocol are included in the Supplementary
Material. \todo{Insert final numerical result and robustness conclusion.}
\end{itemize}

%==============================
\begin{shaded*}
\noindent \textbf{Q:}
{\itshape
Limited evaluation of more complex scenarios: Although the method is
tested with up to 45 movable obstacles, the number of independent
trials is limited.
}
\end{shaded*}

\noindent \textbf{A:}
We agree. The original density sweep used five independent seed sequences per
obstacle density, which is too small for a strong statistical statement at the
highest clutter levels. The revised experiment expands this set to
\todo{insert final number, target at least 10--20} independent seed sequences.
Within one seed, increasing obstacle counts remain nested by adding objects, but
different seeds are independently generated. The revised figure reports
success and timing dispersion rather than means alone.

\noindent\textbf{Changes in the revised manuscript:}
\begin{itemize}
\item Sec.~IV-A6 now states:
\change{The revised evaluation expands the original five seeds to
\todo{insert final number, target at least 10--20} independent seed sequences
per density. Within each seed sequence, denser instances are generated by
incrementally adding objects, while different seeds define independent base
layouts.}

\item The density figure caption now explicitly requires error bars/dispersion:
\change{The revised plot reports success rate and planning/execution time with
dispersion across independent seeds. \todo{Replace the current five-seed figure
by the expanded-seed mean$\pm$standard-deviation/confidence-interval plot.}}
\end{itemize}

%==============================
\begin{shaded*}
\noindent \textbf{Q:}
{\itshape
The experiments primarily use two robots. What happens when the
number of robots increases with the complexity of the task? How does
the search and the contact/force formulation scale to larger robot
teams?
}
\end{shaded*}

\noindent \textbf{A:}
Thank you for raising the team-size question. The revised paper addresses it in
both the complexity analysis and a dedicated simulation study.

The key point is that the implementation does not enumerate all joint contact
tuples. If each robot has $|\mathcal A|$ reachable contact candidates, exhaustive
assignment would require $|\mathcal A|^N$ tuples. Instead, one greedy replacement
sweep tests at most $N|\mathcal A|$ single-contact substitutions. Each test
solves the force LP whose decision vector has dimension $2N$. Thus larger teams
can enlarge the feasible wrench set but also increase contact assignment and LP
cost. The number of strategies entering simulation remains bounded by the
retained-mode and per-expansion candidate caps.

A new table varies $N\in\{1,2,3,4\}$ and reports success, mode-generation time,
planning time, and push count. \todo{Insert final results and interpretation.}

\noindent\textbf{Changes in the revised manuscript:}
\begin{itemize}
\item Sec.~III-E2 now adds:
\change{If $|\mathcal{A}|$ reachable contact candidates are considered per
robot, one greedy replacement sweep tests at most $N|\mathcal{A}|$
single-contact substitutions rather than enumerating $|\mathcal{A}|^N$ joint
tuples. Each substitution solves the LP in Eq.~(10), whose force-vector
dimension is $2N$.}

\item Sec.~IV-A7 and Table~\todo{final table number} provide the dedicated
robot-team scaling experiment. \todo{Insert final numerical data.}
\end{itemize}

%==============================
\begin{shaded*}
\noindent \textbf{Q:}
{\itshape
What guarantees, if any, does the hybrid search provide regarding
solution quality?
}
\end{shaded*}

\noindent \textbf{A:}
We agree that this should have been stated explicitly. No global optimality or
completeness guarantee is intended for the underlying continuous hybrid
problem. PIHS samples a finite set of pushing directions/contact modes, uses
bounded gap lookahead and a heuristic best-first priority, and terminates at the
first physics-validated $W$-clear solution. Moreover, the remaining-cost term is
not required to be admissible. Equation~(5) therefore defines the clearance
cost used to rank schedules rather than an optimization problem that PIHS claims
to solve globally.

\noindent\textbf{Changes in the revised manuscript:}
\begin{itemize}
\item Eq.~(5) is now introduced as the clearance objective
$\mathcal J_{\mathrm{clr}}(\pi)$ rather than as an implicit global-minimization
claim.

\item Sec.~III-D now states:
\change{PIHS is a feasibility-oriented heuristic search. Equation~(5) defines
the schedule cost used to rank candidates, while the implemented search samples
a finite set of pushing directions and modes, uses bounded gap lookahead, and
terminates when a feasible $W$-clear schedule is found. Therefore, no global
optimality or completeness claim is made with respect to the underlying
continuous hybrid planning problem.}

\item The termination statement was revised to say that the returned schedule is
a feasible solution without implying global minimization of
$\mathcal J_{\mathrm{clr}}$.
\end{itemize}

%==============================
\begin{shaded*}
\noindent \textbf{Q:}
{\itshape
The objective includes the total duration T and the control-effort
cost J. Is the time needed for the final vehicle to actually traverse
the cleared corridor included in this objective, or is the objective
only concerned with clearing the obstacles? Similarly, does J only
capture the robots' effort in moving the obstacles?
}
\end{shaded*}

\noindent \textbf{A:}
The objective concerns the robotic clearance operation only. The external
vehicle traversal time is not included in $T$; vehicle passage is a terminal
feasibility condition. Similarly, $J$ represents the robots' manipulation/control
effort during obstacle pushing. We have made both points explicit to remove the
ambiguity.

\noindent\textbf{Changes in the revised manuscript:}
\begin{itemize}
\item Sec.~II-D now states:
\change{$T\triangleq\sum_{k=1}^{K}\Delta t_k$ is the duration of the robotic
clearance operation only; the subsequent traversal time of the external vehicle
is not included. The weight $\alpha>0$ trades time against $J(\cdot)$, which
measures the robots' simulation-estimated control/manipulation effort for moving
obstacles. Vehicle traversal enters the formulation only through the terminal
$W$-clearance condition.}
\end{itemize}

%==============================
\begin{shaded*}
\noindent \textbf{Q:}
{\itshape
The paper demonstrates online replanning when the real execution
differs from the simulation, How can simple pushing robots evaluate
this discrepancy?
}
\end{shaded*}

\noindent \textbf{A:}
Thank you for pointing out that the original wording could be interpreted as
requiring online parameter identification or force sensing. That is not the
implementation. The phrase ``deviates from simulation'' has been replaced by an
explicit state-based execution monitor.

In the hardware experiments, motion capture supplies robot and obstacle poses.
The controller checks whether the pre-contact transition is feasible, whether
robots maintain/track their intended contact poses, whether motion stalls,
whether the task-specific geometric push predicate is achieved, and whether the
updated WCCG still supports the remaining schedule. A failure of any of these
checks triggers replanning from the observed state. Thus simple velocity-controlled
pushers do not need to estimate the underlying mass/friction error itself.

\noindent\textbf{Changes in the revised manuscript:}
\begin{itemize}
\item Sec.~III-E1 now states:
\change{Replanning does not require the robots to estimate an unknown force-model
error directly. The execution monitor uses observed robot/object poses and task
predicates: replanning is triggered when a pre-contact transition is blocked,
contact tracking stalls or loses the planned contact, the pushed object fails to
satisfy the corresponding geometric task predicate within its rollout budget,
or the observed post-push geometry invalidates the remaining schedule/WCCG
connectivity.}

\item The Supplementary Material reports the current contact-position threshold
(0.10~m) and stall check (less than 0.05~m motion over 0.25~s), together with the
task-predicate logic.
\end{itemize}

\newpage
%========================================
\section*{To Reviewer 4}
%========================================
\begin{shaded*}
\noindent \textbf{Q:}
{\itshape
Overall, I think the approach is technically reasonable and the
experiments support the proposed method. My comments are mainly about
clarification rather than major technical concerns, and I would support
acceptance after these revisions.
}
\end{shaded*}

\noindent \textbf{A:}
We sincerely thank the reviewer for the positive assessment and for identifying
four places where the original presentation was insufficiently explicit. The
revision addresses each point directly, as detailed below.

%==============================
\begin{shaded*}
\noindent \textbf{Q:}
{\itshape
First, the WCCG connectivity condition is important since it is used
throughout the search and also determines termination. The intuition is
reasonably clear, but I think the paper would benefit from a little
more justification here, or perhaps a comparison with a standard
configuration-space connectivity check. This seems especially relevant
when the obstacles are non-convex.
}
\end{shaded*}

\noindent \textbf{A:}
We agree. The revised Sec.~III-A now makes the configuration-space interpretation
explicit and adds the suggested independent comparison.

For a vehicle represented by a radius-$W/2$ enclosing disk, inflating obstacle
geometry by $W/2$ converts every pair of boundary components separated by less
than $W$ into overlapping clearance barriers. WCCG bridge--bridge edges encode
these narrow barriers. Non-convex obstacles are first decomposed into
clearance-relevant convex components/boundary segments before the closest-point
bridges are inserted. The same-face condition is therefore interpreted as a
topological query on the planar regions not separated by these inflated barriers.

Because implementation details and degenerate non-convex cases still matter, we
also add an empirical check against a conventional rasterized configuration-space
flood-fill reference, including dedicated non-convex scenes.
\todo{Insert final trial count, resolution, agreement, and false-positive/false-negative results.}

\noindent\textbf{Changes in the revised manuscript:}
\begin{itemize}
\item Sec.~III-A2 now adds:
\change{The geometric interpretation follows the standard configuration-space
inflation for the vehicle-center point. Expanding every obstacle boundary by a
disk of radius $W/2$ transforms a pair of boundary components whose closest
distance is below $W$ into overlapping clearance barriers. The corresponding
WCCG bridge--bridge edge records this non-traversable bottleneck.}

\item The previous statement that WCCG ``eliminates resolution-induced errors''
was softened. The revised remark states that WCCG avoids choosing a workspace
grid resolution but still depends on correct decomposition/bridge construction.

\item Sec.~IV-A3 and Supplementary Sec.~S2 add the requested C-space comparison.
\todo{Insert final result.}
\end{itemize}

%==============================
\begin{shaded*}
\noindent \textbf{Q:}
{\itshape
Second, I was a little confused about the relationship between Eq. (5)
and the actual search procedure. Eq. (5) is written as an optimization
problem, while the implementation uses sampled pushing directions,
limited lookahead, and a heuristic best-first search that terminates
once a feasible solution is found. The authors should clarify whether
any optimality or completeness claim is intended. If not, it may be
better to describe Eq. (5) as the objective used to guide the search.
}
\end{shaded*}

\noindent \textbf{A:}
We agree completely with this interpretation and have revised the formulation
accordingly. Eq.~(5) now defines the clearance cost
$\mathcal J_{\mathrm{clr}}(\pi)$ used to guide candidate ranking. PIHS is
explicitly described as a feasibility-oriented heuristic search with finite
sampling, bounded lookahead, non-admissible heuristic ranking, and first-feasible
termination. No global optimality or completeness claim is intended.

\noindent\textbf{Changes in the revised manuscript:}
\begin{itemize}
\item Sec.~II-D replaces the minimization display by the objective definition:
\change{$\mathcal{J}_{\mathrm{clr}}(\pi)\triangleq
T+\alpha\sum_{k=1}^{K}J(m_k,\boldsymbol\xi_k,\mathbf{s}(t_k)).$}

\item Sec.~III-D explicitly states the absence of global optimality/completeness
guarantees and notes that the remaining-cost term is not required to be
admissible.
\end{itemize}

%==============================
\begin{shaded*}
\noindent \textbf{Q:}
{\itshape
Third, it is not completely clear how the optimized contact forces are
used on the hardware. The pushing mode includes both contact points and
forces, but the execution section mainly describes moving the robots to
the contact configuration and then following the desired pushing
direction. It would be useful to briefly explain the low-level
controller and whether these forces are actually commanded or are
mainly used to check feasibility. Related to this, I am also not sure
whether coupled motion between multiple objects, such as in the second
hardware experiment, is explicitly considered during planning or is
mostly handled by the physics simulation.
}
\end{shaded*}

\noindent \textbf{A:}
Thank you for identifying both ambiguities. The optimized forces are primarily
planning-level feasibility variables; they are not sent to the hardware as
force commands. For a candidate contact tuple, the LP determines whether the
robots can generate a wrench compatible with the desired object motion within
the assumed force/friction constraints. The selected hardware mode is then
executed by closed-loop planar velocity control. Specifically, the Mecanum UGVs
receive ROS \texttt{geometry\_msgs/Twist} commands while tracking the assigned
contact poses and a short forward object-motion reference using motion-capture
feedback.

Coupled object motion is represented through simulation. Although a high-level
strategy names one directly targeted obstacle, $\mathsf{EvalSim}$ advances all
robots and all movable bodies. Object--object contacts are therefore resolved by
the physics engine, and the complete resulting scene state is returned to the
search. The adjacent T-shaped object moved in hardware Scenario~2 is exactly
this type of contact-induced motion rather than a separately scripted task.

\noindent\textbf{Changes in the revised manuscript:}
\begin{itemize}
\item Sec.~III-C2 now states:
\change{The optimized force vector is used here to certify and rank wrench
feasibility under the modeled force limits. On the hardware, the robots execute
pose/velocity commands that maintain these contact locations rather than
directly commanding the optimized forces.}

\item Sec.~III-E1 now specifies the hardware interface:
\change{The controller converts the planned contact geometry and desired object
motion into planar robot velocity commands (ROS \texttt{Twist}); motion-capture
poses provide the feedback used to track the moving contact poses.}

\item Sec.~III-C3 now clarifies coupled objects:
\change{The rollout advances all robots and all movable bodies jointly in the
physics engine. Consequently, object--object contacts induced by pushing the
target obstacle are resolved during $\mathsf{EvalSim}$ and their coupled motion
is retained in the full successor state.}
\end{itemize}

%==============================
\begin{shaded*}
\noindent \textbf{Q:}
{\itshape
The experiments are generally convincing. The paper already mentions
that the baselines use the same W-clearance criterion, object
distribution, and execution simulator. However, a few more
implementation details, such as robot reachability checks, timeout
limits, and computation budget, would make the comparison easier to
interpret. I would also suggest adding error bars or standard
deviations, particularly for the scalability experiment where only five
seeds are used for each obstacle density.
}
\end{shaded*}

\noindent \textbf{A:}
We agree and have expanded both the main setup and the Supplementary protocol.

\begin{itemize}
\item \textbf{Reachability.} A contact candidate is retained only when the
associated pre-contact pose is collision free and the transition planner can
reach it from the current robot configuration. A multi-robot tuple is rejected
if any assigned robot fails this check.

\item \textbf{Common budgets.} The main 40-trial comparison now includes the
shared timeout, search/control cap, and parallel-worker count.
\todo{Verify and insert the exact Scenario-1 values.}

\item \textbf{Density-suite budgets.} The stored evaluation data specify 16
parallel workers, an 800~s per-trial timeout, and a 20,000-control-step cap. The
original manuscript's 400~s statement is corrected to 800~s. Failed and
pruned trials are assigned this same cap in the planning-time statistic.

\item \textbf{Statistics.} The density sweep is expanded from five seeds to
\todo{insert final count}, and the regenerated figure reports variability
(mean$\pm$standard deviation or confidence interval) rather than means alone.
\end{itemize}

\noindent\textbf{Changes in the revised manuscript:}
\begin{itemize}
\item Sec.~III-C2 adds the explicit pre-contact reachability criterion.

\item Sec.~IV-A1 now states:
\change{All compared methods use the same robot footprints, obstacle geometry,
collision/reachability checks, success criterion, and execution simulator. A
candidate that requires an unreachable or collision-infeasible pre-contact
robot configuration is rejected before pushing.}

\item Sec.~IV-A6 now states the density-suite computation budget and replaces
the original 400~s cap by the correct 800~s timeout. The figure caption also
requires the expanded-seed error-bar plot.
\end{itemize}

%==============================
\end{document}
